\question{نمونه‌ی کدنویسی}

این فایل نشان می‌دهد چطور کد را در پاسخ‌ها قرار دهید.
برای نوشتن کلمات کلیدی کوتاه داخل متن، می‌توانید از \code{...} استفاده کنید؛ مثلاً \code{for} یا \code{while}.

\subsection{کد کوتاه}
قالب به‌صورت پیش‌فرض بدون نیاز به \lr{-shell-escape} کامپایل می‌شود (حالت \lr{listings}).
اگر خواستید خروجی بهتر \lr{minted} داشته باشید، در \lr{\texttt{main.tex}} بنویسید:
\lr{\texttt{\textbackslash usepackage[minted]\{commons/style\}}}

\subsection{یک قطعه کد}
\begin{listing}[H]
	\begin{latin}
		\begin{minted}{python}
def gcd(a, b):
    while b:
        a, b = b, a % b
    return a
		\end{minted}
	\end{latin}
	\caption{نمونه‌ی کد (Python)}
\end{listing}

\subsection{یک نکته‌ی کوچک}
اگر می‌خواهید داخل پاسخ، نام فایل/دستور را واضح بنویسید:
\begin{ابجد}
	\item فایل: \lr{\texttt{answers/template\_code.tex}}
	\item دستور: \lr{\texttt{\textbackslash lr\{...\}}}
\end{ابجد}

\subsection{نکته‌ی متنی}
داخل متن فارسی هم می‌توانید عبارت انگلیسی بیاورید: \lr{time complexity}، یا حتی نام تابع مثل \lr{gcd}.

